\project{19}{SPI LCD tilt meter}{Raspberry Pi 3 + Waveshare 3.5 inch LCD (A) + SEN0032 ADXL345}{the LCD as an offline field instrument}

\begin{unique}
Owns the 3.5 inch LCD as a field instrument with ADXL345 tilt. No broker, no radio. The LCD owns the 40-pin header, so this cannot run with P06 on the same Pi at the same time, a different sitting.
\end{unique}

\subsection*{Intent}

A Pi 3, a Waveshare 3.5 inch (A) panel and an ADXL345. Pitch and roll on the glass, drawn as two bars, and nothing else. The instrument is battery-unaware: it is an offline meter, not a gateway, and there is no part of it that cares whether a network exists.

That absence is the lab. Every other display in this book is attached to something: P12 shows broker topics, P06 shows ranges arriving from a Nucleo, P04 shows a waterfall from a USB probe. This one shows a quantity it measured itself, on a panel it drives itself, with no client and no uplink. If an MQTT client is running during the acceptance test, the lab has drifted into being a different lab.

\begin{kit}
Pi 3, Waveshare 3.5 inch RPi LCD (A), SEN0032 ADXL345, jumpers for the sensor. This panel has no pass-through, so read the wiring note before planning the sitting.
\end{kit}

\begin{rulebox}{The LCD owns the 40-pin header}
This lab and P06 both want the 3.5 inch LCD on a Pi 3. They cannot share one Pi in the same sitting. Build one, accept it, tear it down, then build the other.
\end{rulebox}

\subsection*{System architecture}

\diagram{p19_arch}{An instrument with no network side. SPI carries the panel, I2C carries the accelerometer, and the loop between them never leaves the board.}

There are exactly two buses in this lab and both of them stay on the host. SPI drives the panel through the vendor overlay, which on Bookworm means the current Waveshare note for the 3.5 inch (A) rather than the flexfb recipes that circulated in 2016. I2C1 carries the ADXL345 at 0x53, and reaching it is the difficulty of this lab rather than a detail of it, because the panel covers the pins the bus lives on. Between them sits one Python loop that reads six registers, computes two angles and paints two bars.

Nothing else is present. There is no broker, no radio, no USB gadget and no second host. Confirm that a framebuffer exists before writing any drawing code, exactly as P06 does: either \texttt{/dev/fb0} or a DRM card, depending on what the vendor overlay gives you on your image.

\begin{note}{This lab and P06 use the same panel}
P06 puts the same Waveshare 3.5 inch (A) on a Pi 3 and paints an 8x8 range grid that arrives over USB from a Nucleo. This lab paints two bars from a sensor on the host's own I2C bus. They share the panel, the host and the header, so they cannot share a sitting. The overlay work is the same in both, which is why this chapter points at P06 for it rather than repeating it.
\end{note}

\subsection*{Wiring and schematic}

\diagram{p19_schematic}{The panel covers the pins the sensor needs, and the two ways past that. SDO left floating selects 0x53, the same address the sensor uses in P01 and P09.}

\begin{table}[H]\centering\small
\begin{tabular}{p{42mm}p{26mm}p{22mm}p{58mm}}
\toprule
Signal & Pi header pin & BCM GPIO & Note \\
\midrule
ADXL345 VCC & 1 & 3V3 & Never 5 V \\
ADXL345 SDA & 3 & GPIO2 & I2C1 data \\
ADXL345 SCL & 5 & GPIO3 & I2C1 clock \\
ADXL345 GND & 6 & GND & \\
ADXL345 SDO & & & Left floating, selects 0x53 \\
LCD panel & 40-pin header & & The HAT owns the header for the sitting \\
\bottomrule
\end{tabular}
\caption{Wiring as the source gives it, which this panel cannot be wired to. The four rows above assume a pass-through, and the Waveshare 3.5 inch (A) has none. Read the note before touching a jumper.}
\end{table}

\begin{note}{This panel has no pass-through, and the lab is blocked until that is worked around}
The source says to wire the sensor to pins 1, 3, 5 and 6 "if the LCD HAT pass-through exposes them (it usually does)". On this panel it does not. The board mates with the first 26 pins through a female header on its underside and carries no header on top, so that whole corner sits under the panel and no jumper reaches it.

The arithmetic is worse than needing different pins. Hardware I2C1 exists only on pins 3 and 5, so there is no second I2C bus to move to. Pin 1 and pin 17 are the only 3.3 V pins on the header and both are inside the covered range, so with the panel seated there is no supply pin left either. What remains above pin 26 is grounds on 30, 34 and 39 and GPIOs on 29, 31, 32, 33, 35, 36, 37, 38 and 40, with pins 27 and 28 left alone as the identity EEPROM lines.

One detail from the panel's own pinout is what makes a fix possible rather than impossible: it declares pins 3 and 5 not connected. The I2C bus is electrically free and only physically buried, so nothing is contending for it.

\textbf{There was a second way out, and it is closed.} The attractive route was to leave the panel seated, take 3.3 V from the POW-BB rail that P15 already labels, ground to pin 30, 34 or 39, and run a bit-banged bus through the \texttt{i2c-gpio} overlay on two of the free GPIOs above pin 26. It was attractive because it costs nothing and because it would have made P15's rails load-bearing in a later lab rather than a one-off.

It does not work, and the carrier's own schematic is what says so. Its Raspberry Pi connector symbol is drawn as thirteen rows of two, numbered 1 and 2 down to 25 and 26, and it stops there; the netlist carries twenty-six pads and nothing above. Pins 27 to 40 are not contacts on this board and are not listed even as not connected. The carrier outline is 85.06 by 56.21 mm against the Pi's 85 by 56 mm, so the board runs the full length of the header and past it, and pin 1 of the socket sits at the SD card and power end. Those fourteen pins are electrically untouched and mechanically buried the moment the panel is seated, so there is nowhere to attach a jumper to them.

\textbf{So the lab needs one part: a 2x20 stacking header, or a tall header whose pins 27 to 40 clear the 26-pin socket.} It is not on this bench. That is a better position than the lab held before, because it is no longer blocked on something unknown; it is blocked on a purchase it can name.

Settled on Thursday 8 October 2026, and worth recording how. An earlier version of this note asked for a ruler along the carrier's long edge, and an inference from the vendor's wiki sentence and the glass panel's dimension drawing already pointed at this same answer. That inference was deliberately not written down, because it rested on a marketing phrase and a drawing of a different part. The schematic settled it without any measuring at all, which is the lesson worth keeping: the cheapest measurement is often a document nobody had opened.
\end{note}

\begin{asciiart}
ADXL345                 Pi 40-pin (NOT reachable: the panel covers pins 1-26)
VCC ------------------- pin 1  3V3      NEVER 5V
GND ------------------- pin 6  GND
SDA ------------------- pin 3  GPIO2
SCL ------------------- pin 5  GPIO3
SDO floating => 0x53

Waveshare 3.5" LCD (A) seated on the Pi 3 40-pin header: this lab's HAT.
Use the current Bookworm vendor overlay, not the 2016 flexfb recipe.
This cannot share a Pi with P06 in the same sitting.
\end{asciiart}

\subsection*{Bench layout}

\diagram{p19_bench}{Bench layout. The board is meant to be picked up and tilted, so the sensor has to be fixed to the same body as the panel and the jumpers have to be short enough to survive being moved.}

The instrument gets handled during its own acceptance test, which is unusual for this book. Fix the ADXL345 to the same piece of board or case as the Pi, and keep the four jumpers short. A sensor that swings on long leads while you rotate the assembly reports its own swing as tilt, and the bars will tell you so.

Note the sensor's orientation on the assembly before the first rotation, because pitch and roll are defined by which axis points where. There is no calibration step in this lab, so the mapping from a physical rotation to a bar is whatever the mounting made it. If the wrong bar moves, rotate the sensor rather than editing the formulas.

\subsection*{Software design (UML)}

\diagram{p19_uml}{The activity view of one pass: read the six registers, convert to g, compute pitch and roll, draw two bars, repeat. There is no branch in this diagram that reaches a network.}

The loop is deliberately flat. Read six bytes starting at 0x32, convert each axis pair to g the way P01 does, compute the two angles, and draw. Pitch is $\arctan$ of $a_x$ against $\sqrt{a_y^{2} + a_z^{2}}$, and roll is $\arctan$ of $a_y$ against $\sqrt{a_x^{2} + a_z^{2}}$, both taken with the two-argument form so the sign survives. Everything after that is drawing: two bars, full scale at the ends of the range you chose.

One thing is worth stating because it is easy to want: there is no filter here and no calibration step. The source gives the two formulas and the drawing, and that is the whole product. Anything more belongs to the labs that own the physics of orientation, not to a panel demonstration.

The other thing the diagram shows by omission is that there is no failure path worth drawing. An instrument with a network has states for connecting, retrying and giving up. This one has a read that either returns six bytes or raises, and a draw that either paints or does not. That is the trade the chapter is making: less reach, far less to go wrong.

\subsection*{Data flow (ASCII)}

\begin{asciiart}
  Pi 3, offline
  +------------------------------------------------------------------+
  |  i2c-1 -> 0x53 ADXL345                                            |
  |     |  read 6 bytes at 0x32                                       |
  |     v                                                             |
  |  ax, ay, az  (g, full resolution)                                 |
  |     |                                                             |
  |     +--> pitch = atan2(ax, sqrt(ay*ay + az*az))                   |
  |     +--> roll  = atan2(ay, sqrt(ax*ax + az*az))                   |
  |               |                                                   |
  |               v                                                   |
  |  two bars on the framebuffer  ->  SPI  ->  3.5" LCD (A), 480x320  |
  +------------------------------------------------------------------+

  no broker, no radio, no USB gadget, no second host
  no MQTT client is running during the acceptance test
\end{asciiart}

\subsection*{Steps}

\begin{steps}
\item \textbf{Seat the LCD} on a powered-off Pi 3 and confirm the vendor overlay as in P06. On Bookworm that means the current Waveshare 3.5 inch (A) note, not the flexfb recipe from 2016. Confirm that a framebuffer exists before writing any drawing code.
\begin{shellcode}
ls /dev/fb* /dev/dri/card* 2>/dev/null
\end{shellcode}

\item \textbf{Wire the ADXL345} by whichever of the two routes in the wiring note applies to the panel in front of you, because the printed pins 1, 3, 5 and 6 are underneath it. Then confirm the address.
\begin{shellcode}
sudo apt-get install -y i2c-tools python3-smbus
sudo i2cdetect -y 1          # must show 53
\end{shellcode}

\item \textbf{Read the sensor and compute the two angles.} The register sequence is the one from P01: 0x2D to enter measure mode, 0x31 for full resolution, then six bytes from 0x32.
\begin{pycode}
import math, time, smbus
bus = smbus.SMBus(1); A = 0x53
bus.write_byte_data(A, 0x2D, 0x08)   # measure
bus.write_byte_data(A, 0x31, 0x0B)   # full-res +/-16g

def g(lo, hi):
    v = (hi << 8) | lo
    if v & 0x8000:
        v -= 65536
    return v / 256.0

while True:
    b = bus.read_i2c_block_data(A, 0x32, 6)
    ax, ay, az = g(b[0], b[1]), g(b[2], b[3]), g(b[4], b[5])
    pitch = math.degrees(math.atan2(ax, math.sqrt(ay*ay + az*az)))
    roll  = math.degrees(math.atan2(ay, math.sqrt(ax*ax + az*az)))
    draw_two_bars(pitch, roll)       # on the framebuffer
    time.sleep(0.05)
\end{pycode}

\item \textbf{Draw two bars on the framebuffer.} One bar for pitch, one for roll, each with a marked centre and marked ends. The panel is 480x320, so a horizontal bar per angle across most of the width leaves room for a numeric readout if you want one.

\item \textbf{Check that nothing else is running.} No MQTT client, no publisher, no uplink of any kind. That is part of the acceptance test rather than a tidiness preference.
\begin{shellcode}
pgrep -a mosquitto_sub mosquitto_pub || echo "nothing publishing, good"
\end{shellcode}

\item \textbf{Tear the LCD down} when the lab is accepted, so the Pi 3 is free for P06 in a later sitting.
\end{steps}

\subsection*{Acceptance test}

\begin{itemize}
\item Rotate the assembly 90\textdegree{} on one axis: that bar travels full scale.
\item The other bar stays near zero during the same rotation.
\item No MQTT client is running.
\end{itemize}

\subsection*{Practices}

Keep the instrument offline for the whole sitting. The value of this lab is that a display and a sensor can be a complete product with no network behind them, and adding an uplink to see the numbers somewhere else removes the only thing this chapter owns.

One 40-pin HAT per host, as everywhere in this book. The 3.5 inch LCD is the HAT here, so the Pi 3 cannot also carry the MCC 118, a cellular HAT or the Explorer700, and it cannot be running P06 at the same time. Power off before unseating the panel, and put the ADXL345 back in the bin with its jumpers rather than leaving it dangling from a stored board.

The sensor is the same SEN0032 used in P01 and P09, and it is the same address and the same register sequence in all three. When it appears at 0x53 here, that is a part behaving as it has behaved twice already, and when it does not, the fault is in how the bus was reached rather than in the chip.

\subsection*{Pitfalls}

\begin{itemize}
\item \textbf{Following a 2016 flexfb recipe.} On Bookworm the Waveshare 3.5 inch (A) uses the current vendor overlay. The old recipe produces a panel that stays dark and a stack of changes to undo.
\item \textbf{Writing drawing code before checking for a framebuffer.} Confirm \texttt{/dev/fb0} or a DRM card first, exactly as P06 does.
\item \textbf{Long jumpers to the sensor.} The assembly gets rotated during its own test. A sensor swinging on loose leads measures its own swing.
\item \textbf{5 V on the accelerometer.} The SEN0032 goes to pin 1, never pin 2.
\item \textbf{Assuming a pass-through.} The source says one usually exists, and this panel has none, which is what blocked the lab. Read the board's own pinout before planning any wiring that depends on reaching a pin underneath it.
\item \textbf{Trying to run P06 on the same Pi in the same sitting.} Both labs want the LCD on the 40-pin header. Tear one down first.
\item \textbf{Leaving a subscriber running from an earlier lab.} The acceptance test says no MQTT client is running, and a forgotten \texttt{mosquitto\_sub} in another terminal is still a client.
\item \textbf{Using the single-argument arctangent.} The sign of the angle is then lost across half the range, and one bar travels the wrong way.
\end{itemize}

\subsection*{Sources}

\begin{itemize}
\item Waveshare 3.5 inch RPi LCD (A) wiki, and the current note for Bookworm, \url{https://www.waveshare.com/wiki/3.5inch_RPi_LCD_(A)}
\item DFRobot SEN0032 ADXL345 product wiki, \url{https://wiki.dfrobot.com/}
\item Analog Devices ADXL345 data sheet, register map 0x2D, 0x31, 0x32 and the SDO address strap
\item Raspberry Pi device tree and overlay documentation, \url{https://www.raspberrypi.com/documentation/computers/configuration.html}
\item Linux framebuffer and DRM documentation, for the difference between \texttt{/dev/fb0} and a DRM card
\item \texttt{i2c-tools} manual pages, \texttt{i2cdetect(8)}
\item Python \texttt{math} module documentation, \texttt{atan2} and its sign convention
\end{itemize}
